\chapter{Using and Creating APIs}

As a software engineer, you are likely familiar with building
applications that interact with various external services and data
sources. One of the most common methods for communication and
integration is through HTTP APIs (Application Programming
Interfaces). HTTP APIs provide a standardized way for applications to
exchange data and functionality over the internet.\index{Web APIs} \index{HTTP}

This chapter will introduce you to the world of HTTP APIs and explore
how you can leverage them in your software development projects. We
will cover the fundamental concepts, techniques, and best practices
for effectively working with HTTP APIs.

An HTTP API allows two software systems to communicate and exchange
information using the Hypertext Transfer Protocol (HTTP). It enables
your application to make requests to an API server and receive
responses in a structured format, such as JSON (JavaScript Object
Notation) or XML (eXtensible Markup Language).

Using HTTP APIs offers a range of benefits for software engineers. It
allows you to leverage external services and data sources, enabling
your application to access functionality or retrieve valuable
information from third-party systems. This opens up opportunities for
integration with popular platforms, social media networks, payment
gateways, geolocation services, and much more.

Throughout this chapter, we will explore various aspects of working
with HTTP APIs, including:

\begin{itemize}
  \item API endpoints and methods: Understanding how to interact with an
        API involves identifying the available endpoints and the supported
        methods, such as GET, POST, PUT, DELETE, and so on. We will discuss how to
        construct API requests and handle different response formats.

  \item Authentication and authorization: Many APIs require
        authentication to ensure secure access and protect sensitive
        data. We will delve into different authentication mechanisms,
        including API keys, tokens, OAuth, and other authentication
        protocols commonly used in API integrations.

  \item Request parameters and payloads: APIs often accept additional
        parameters or payloads to customize the request or send data for
        processing. We will explore how to pass query parameters, request
        headers, and request bodies when interacting with APIs.

  \item Error handling and status codes: Learning how to handle errors
        and interpret status codes returned by APIs is crucial for building
        robust and resilient applications. We will discuss common status
        codes and best practices for handling various scenarios gracefully.

  \item Rate limiting and throttling: Many APIs impose restrictions on
        the number of requests you can make within a given timeframe to
        prevent abuse and ensure fair usage. We will cover techniques for
        handling rate limiting and implementing efficient strategies to
        manage API quotas.

  \item API documentation and testing: Proper documentation is essential
        for understanding an API's capabilities, endpoints, and expected
        behavior. We will explore how to read and interpret API
        documentation, as well as techniques for testing and validating API
        integrations.
\end{itemize}

By mastering the art of using HTTP APIs, you will expand your
development toolkit and gain the ability to seamlessly integrate your
applications with external services, leverage their functionalities,
and build powerful, interconnected systems.

So, let's dive into the world of HTTP APIs and uncover the endless
possibilities they offer for enhancing your software engineering
projects.

% Written by Arjan June-July 2026
\section{Setting up an API}

Create a file called \texttt{server.js}. This will be a basic webserver that we will use to test our an API. This is called a \newterm{RESTful API}. Write the following code to create a basic server for our API.
\begin{minted}{javascript}
const express = require("express");

const app = express();
const PORT = 3000;

const student = {
  student_id: 42,
  name: "John Smith",
  major: "Computer Science",
  credits_completed: 88
};

app.get("/students/42", (req, res) => {
  res.json(student);
});

app.listen(PORT, () => {
  console.log(`API running at http://localhost:${PORT}`);
});
\end{minted}
Then, run the following
\begin{minted}{shell}
npm init -y   
npm install express
node server.js
\end{minted}
This will create a new server, along with a ton of dependency files that javascript and 
node automatically install as part of your server.
\subsection{What is an endpoint?}
An \newterm{endpoint} is a URL exposed by the API. It usually is able to display a page or 
data based on the request of the endpoint.

We can navigate to the endpoint \texttt{GET /students/42}, where students is the resource 
that usually maps to a database, and \texttt{42} is the id of a student in the database. 
To retrieve this student from the database, we can navigate to
\begin{minted}{python}
GET http://localhost:3000/students/42
\end{minted}
which is an endpoint from our API. This will output
\begin{minted}{json}
{
  "student_id": 42,
  "name": "John Smith",
  "major": "Computer Science",
  "credits_completed": 88
}
\end{minted}

Note that this is done in json format, common for backend data parsing.

Each endpoint combines a resource path with an HTTP method to complete an action.

This leads us to the idea of \newterm{CRUD operations}. CRUD, which stands for CREATE, 
REMOVE, UPDATE, DELETE, are a form of outlining an API's routes. Calling a route that 
aims to alter should follow one of the CRUD operations. GET does not have an operation 
since it is \newterm{read-only}.

\begin{table}[]
  \begin{tabular}{|l|l|l|l|}
    \hline
    \textbf{HTTP Method} & \textbf{Purpose}       & \textbf{CRUD Op} & \textbf{Example}                  \\ \hline
    GET                  & Fetch Data             & N/A              & Get student information           \\ \hline
    PUT                  & Update Data            & UPDATE           & Fully update student              \\ \hline
    PATCH                & Update section of Data & UPDATE           & Update a student's major          \\ \hline
    POST                 & Create Data            & CREATE           & Create a new student              \\ \hline
    DELETE               & Remove Data            & DELETE           & Remove a student after graduation \\ \hline
  \end{tabular}
\end{table}

For a database representing students (assume more than just one), we would want to have 
routes that can complete each CRUD operation.

Before an API can receive data from a client, Express must parse JSON
request bodies. Add the following line after creating the application:

\begin{minted}{javascript}
app.use(express.json());
\end{minted}

This middleware automatically converts incoming JSON into a JavaScript
object that can be accessed through \texttt{req.body}.

\subsection{Reading all resources}

Most APIs provide a route for retrieving an entire collection. Change the students array 
to look like the following.
\begin{minted}{javascript}
const students = [
  {
    student_id: 42,
    name: "John Smith",
    major: "Computer Science",
    credits_completed: 88
  },
  {
    student_id: 43,
    name: "Jane Doe",
    major: "Mathematics",
    credits_completed: 61
  }
];

app.get("/students", (req, res) => {
  res.json(students);
});
\end{minted}

The endpoint

\begin{minted}{text}
GET /students
\end{minted}

returns every student currently stored in the collection. Notice that we did not include an id. 
Now, there are more than one possible ids, so we need to make a general route that takes in an id 
of a student to fetch
\begin{minted}{javascript}
app.get("/students/:id", (req, res) => {
  // use the id that is passed in the route
  const id = Number(req.params.id);
  // find a student from our students array
  const student = students.find((s) => s.student_id === id);

  // if the id is not found, we assume the student is not in the array
  if (!student) {
    return res.status(404).json({ error: "Student not found" });
  }

  res.json(student);
});
\end{minted}
Start the server. Now, \mintinline{bash}|http://localhost:3000/students| returns the 
entire students array, while \mintinline{bash}|http://localhost:3000/students/43| only returns 
information about Jane Doe.

Routing to \mintinline{bash}|http://localhost:3000/students/19| is invalid, so it returns
\begin{minted}{json}
{"error":"Student not found"}
\end{minted}

\section{Requests with cURL}
Curl, sometimes stylized as \texttt{cURL} or \texttt{curl}, is a command-line based resource 
used for sending API and HTTP requests for web server data.

The simplest curl command
\begin{minted}{shell}
curl http://example.com
\end{minted}
which will return the \texttt{html} source code of the website \url{http://example.com}.
\begin{minted}{html}
<!doctype html>
<html lang="en">

<head>
    <title>Example Domain</title>
    <link rel="icon" href="data:,">
    <meta name="viewport" content="width=device-width, initial-scale=1">
    <style>
        body {
            background: #eee;
            width: 60vw;
            margin: 15vh auto;
            font-family: system-ui, sans-serif
        }

        h1 {
            font-size: 1.5em
        }

        div {
            opacity: 0.8
        }

        a:link,
        a:visited {
            color: #348
        }
    </style>
</head>

<body>
    <div>
        <h1>Example Domain</h1>
        <p>This domain is for use in documentation examples without needing permission. Avoid use in operations.</p>
        <p><a href="https://iana.org/domains/example">Learn more</a></p>
    </div>
</body>

</html>
\end{minted}
\section{Updating and Deleting Resources with cURL}

So far our API can read students and create new ones. To finish out
the CRUD table from earlier, we need routes for \texttt{PUT},
\texttt{PATCH}, and \texttt{DELETE}, and we need to know how to
trigger each of them with curl.

\subsection{Fully replacing a resource with PUT}

A \texttt{PUT} request means \emph{replace this entire resource with
  what I'm sending you.} If you leave out a field, that field is
expected to disappear or reset, because the whole object is being
overwritten.

\begin{minted}{javascript}
app.put("/students/:id", (req, res) => {
  const id = Number(req.params.id);
  const index = students.findIndex((s) => s.student_id === id);
 
  if (index === -1) {
    return res.status(404).json({ error: "Student not found" });
  }
 
  // PUT replaces the whole object, so we overwrite it completely
  students[index] = { 
    student_id: id, ...req.body // this format keeps the same student_id and updates the rest of the fields with the new data
  }; 
  res.json(students[index]);
});
\end{minted}

\begin{minted}{shell}
curl -X PUT http://localhost:3000/students/43 \
  -H "Content-Type: application/json" \
  -d '{
      "name": "Jane Doe",
      "major": "Physics",
      "credits_completed": 64
    }'
\end{minted}

Notice we did not send \texttt{student\_id} in the body -- the route
pulls it from the URL instead. Try leaving out
\texttt{credits\_completed} entirely and re-running the request. Because
PUT replaces the whole object, that field will not be included in the
response. This is the behavior that distinguishes PUT from PATCH.

\subsection{Partially updating a resource with PATCH}

A \texttt{PATCH} request means \emph{only change the fields I'm
  sending you, and keep everything else the same.} This is almost always
what you want when, say, updating a single field like a student's
major.

\begin{minted}{javascript}
app.patch("/students/:id", (req, res) => {
  const id = Number(req.params.id);
  const student = students.find((s) => s.student_id === id);
 
  if (!student) {
    return res.status(404).json({ error: "Student not found" });
  }
 
  // PATCH merges the incoming fields onto the existing object
  Object.assign(student, req.body);
  res.json(student);
});
\end{minted}

\begin{minted}{shell}
curl -X PATCH http://localhost:3000/students/43 \
  -H "Content-Type: application/json" \
  -d '{ "major": "Physics" }'
\end{minted}

Run a \texttt{GET} on \texttt{/students/43} afterward and confirm that
\texttt{name} and \texttt{credits\_completed} are untouched; only
\texttt{major} changed. That single-field behavior is the whole point
of PATCH.

\subsection{Removing a resource with DELETE}

A \texttt{DELETE} request removes a resource. Unlike POST, PUT, and
PATCH, it typically does not need a request body; the URL alone
tells the server what to remove.

\begin{minted}{javascript}
app.delete("/students/:id", (req, res) => {
  const id = Number(req.params.id);
  const index = students.findIndex((s) => s.student_id === id);
 
  if (index === -1) {
    return res.status(404).json({ error: "Student not found" });
  }
 
  students.splice(index, 1);
  // status 204 means "success, but there is nothing to send back"
  res.status(204).send();
});
\end{minted}

\begin{minted}{shell}
curl -X DELETE http://localhost:3000/students/44
\end{minted}

Run \texttt{GET /students} again and Jerry Doe should be gone. If you
run the same \texttt{DELETE} command a second time, you should get a
\texttt{404}, since the student no longer exists to delete.

\subsection{Headers with \texttt{-i} and \texttt{-v}}

Every curl command so far has only shown you the body of the response.
But remember: every HTTP exchange also includes headers and a status
code, even when curl hides them by default. Two flags make that
conversation visible.

\texttt{-i} includes the response headers above the body:

\begin{minted}{shell}
curl -i http://localhost:3000/students/43
\end{minted}

\begin{minted}{text}
HTTP/1.1 200 OK
Content-Type: application/json; charset=utf-8
Content-Length: 78
 
{"student_id":43,"name":"Jane Doe","major":"Physics","credits_completed":64}
\end{minted}

\texttt{-v} (verbose) goes further and shows the \emph{request} curl
sent, not just the response it received -- useful when a request is
failing and you aren't sure what curl is actually sending:

\begin{minted}{shell}
curl -v -X PATCH http://localhost:3000/students/43 \
  -H "Content-Type: application/json" \
  -d '{ "major": "Physics" }'
\end{minted}

Reading a \texttt{-v} output line by line is one of the best habits a
beginner can build: lines starting with \texttt{>} are what curl sent,
lines starting with \texttt{<} are what the server sent back.

\section{Authentication and API Keys}
\index{API keys}\index{Authorization header}

Every route we've built so far is completely open: anyone who knows
the URL can read, create, update, or delete students. Real APIs
usually can't work that way, since that data might be private, or
changing it might cost money (imagine n API that sends text messages
or charges a credit card). To restrict access, APIs commonly require
an \newterm{API key}: a long, unique string that identifies who is
making the request.

An API key is typically sent as a request header, most often
\texttt{Authorization} or a custom header like \texttt{X-API-Key}:

\begin{minted}{shell}
curl http://localhost:3000/students \
  -H "Authorization: Bearer abc123yourkeyhere"
\end{minted}

On the server, you check for that header before doing anything else:

\begin{minted}{javascript}
const API_KEY = "abc123yourkeyhere";
 
function requireApiKey(req, res, next) {
  const authHeader = req.headers.authorization;
 
  if (authHeader !== `Bearer ${API_KEY}`) {
    return res.status(401).json({ error: "Unauthorized" });
  }
 
  next();
}
 
app.get("/students", requireApiKey, (req, res) => {
  res.json(students);
});
\end{minted}

If the header is missing or wrong, the client gets back a
\texttt{401 Unauthorized} instead of the data.

A few rules are worth internalizing early, because they matter far
beyond this textbook:

\begin{itemize}
  \item \textbf{Never commit an API key to a public repository.} Once a
        key is pushed to code repository such as GitHub, assume it is compromised, even if you
        delete it in a later commit. If you do accidentally commit a key, revoke it immediately and generate a new one.
  \item \textbf{Never hardcode a real key directly in your source code.} The example does for simplicity. Real projects store keys in environment variables (often loaded from a \texttt{.env}
        file that is excluded from version control or public repositories) and read them with
        something like \texttt{process.env.API\_KEY}.
  \item \textbf{Treat an API key like a password.} Anyone who has your
        key can act as you. They could be reading your data, spending your quota, or
        in some APIs, spending your money.
  \item \textbf{You can only access API's with their specific key.} If you have multiple API's, each one will have its own key. You cannot use the same key for different API's. For example, Spotify's public API has its own key, and a method for you to generate your own key. You cannot use a key from another API (such as OpenAI) to access Spotify's API. The strength here is that if one key is compromised, it does not compromise all of your API's. Secondly, this way, if a breach of data is attempted, Spotify knows who owns the key and can revoke it. Feel free to learn more about Spotify's API and how to generate your own key at \url{https://developer.spotify.com/documentation/web-api/}. Becoming increasingly more common is OpenAI's API, which has its own key and method for generating your own key. You can learn more about OpenAI's API at \url{https://platform.openai.com/docs/guides/authentication}.
\end{itemize}

This is the difference between \newterm{authentication} (proving
\emph{who} you are, usually with a key or token) and
\newterm{authorization} (determining \emph{what} you're allowed to do
once you've proven who you are). A key can identify you correctly and
still be denied access to a specific route if you don't have
permission for it.

\section{Creating resources with POST}

That's all great, but we can put cURL to work on much more laborious tasks. Going back to our created API example, we can post, update, and delete resources using cURL.

A POST request creates a new resource. Unlike GET requests, POST
requests send data inside the request body.

\begin{minted}{javascript}
app.post("/students", (req, res) => {
  const student = req.body;

  // basic validation: ensure a body is received
  if (!student) {
    return res.status(400).json({ error: "Invalid student data" });
  }

  students.push(student);

  // status 201 means that the student was pushed successfully
  res.status(201).json(student);
});
\end{minted}

Notice the use of

\begin{minted}{javascript}
req.body
\end{minted}

which contains the JSON sent by the client.

HTTP requests are composed of headers and a body. Headers are metadata such as \texttt{Content-Type} (most commonly application/json) or \texttt{Authorization} that describe how the request should be handled. The body is the payload that carries actual data, like the student record being created by a POST request.

In Express, request headers are available on \texttt{req.headers}. For JSON payloads, middleware like \texttt{express.json()} parses the body so that the data is available on \texttt{req.body}.

A template for sending json in the body looks something like the following:
\begin{minted}{bash}
curl -X METHOD URL \
     -H "Header: Value" \
     -d 'JSON DATA'
\end{minted}

To send a request body, include the body data in the HTTP request and set the appropriate \texttt{Content-Type} header. For example, with curl:

\begin{minted}{shell}
curl -X POST http://localhost:3000/students \
  -H "Content-Type: application/json" \
  -d '{
      "student_id": 44,
      "name": "Jerry Doe",
      "major": "Computer Science",
      "credits_completed": 23
    }'
\end{minted}
This can be broken down as:

\begin{itemize}
  \item \texttt{curl}: The command-line tool for making HTTP requests.
  \item \texttt{-X POST}: Use the POST HTTP method.
  \item \texttt{http://localhost:3000/students}: The endpoint being contacted. Must exist in our server routes.
  \item \texttt{-H "Content-Type: application/json"}: Tell the server the request body contains JSON.
  \item \verb|-d '{ ... }'|: Send this JSON as the request body.
\end{itemize}

The request body in this example is the JSON string passed with \texttt{-d}. Express will parse that JSON into \texttt{req.body} when \texttt{express.json()} middleware is enabled.

Note that our server.js students array will not update with Jerry Doe, because the server only stores the students array upon initialization, but we can see our now three students by running a GET request on \texttt{/students}

\begin{minted}{json}
[
  {
    "student_id": 42,
    "name": "John Smith",
    "major": "Computer Science",
    "credits_completed": 88
  },
  {
    "student_id": 43,
    "name": "Jane Doe",
    "major": "Mathematics",
    "credits_completed": 61
  },
  {
    "student_id": 44,
    "name": "Jerry Doe",
    "major": "Computer Science",
    "credits_completed": 23
  }
]
\end{minted}


\section{Postman}
\index{Postman}
 
Typing curl commands by hand is a great way to \emph{learn} how HTTP
requests actually work, because every piece of the request (
the method, the headers, the body) has to be spelled out yourself.
But once you already understand what those pieces are, typing them out
every single time gets tedious and error-prone, especially for requests
with long JSON bodies or several headers.
 
\newterm{Postman} is a GUI application for building, sending, and
saving HTTP requests. Everything you've been typing as curl flags has
a matching box in Postman:
 
\begin{table}[]
  \begin{tabular}{|l|l|}
    \hline
    \textbf{cURL}                         & \textbf{Postman}                 \\ \hline
    The URL itself                        & The URL bar                      \\ \hline
    \texttt{-X METHOD}                    & The method dropdown               \\ \hline
    \texttt{-H "Header: Value"}           & The Headers tab                   \\ \hline
    \texttt{-d '\{...\}'}                 & The Body tab (raw / JSON)         \\ \hline
    \texttt{-i} / \texttt{-v} output      & The Response panel and Console    \\ \hline
  \end{tabular}
\end{table}
 
\subsection{Building your first request in Postman}
 
\begin{enumerate}
  \item Open Postman (download available here \href{https://www.postman.com/}) and click \textbf{New} $\rightarrow$
        \textbf{HTTP Request}.
  \item Set the method dropdown (left of the URL bar) to \texttt{GET},
        and enter \texttt{http://localhost:3000/students} in the URL bar.
  \item Click \textbf{Send}. The response body appears in the panel
        below, along with the status code and response time: the same
        information \texttt{-i} shows you in curl, just laid out visually
        instead of printed as text.
\end{enumerate}
 
Now recreate the POST request from earlier:
 
\begin{enumerate}
  \item Change the method to \texttt{POST} and set the URL to
        \texttt{http://localhost:3000/students}.
  \item Open the \textbf{Body} tab, select \textbf{raw}, and choose
        \textbf{JSON} from the format dropdown on the right. Postman will
        automatically set the \texttt{Content-Type: application/json}
        header for you. This is one thing it does that curl makes you
        do by hand.
  \item Paste in a student object:
        \begin{minted}{json}
{
  "student_id": 45,
  "name": "Alex Rivera",
  "major": "Biology",
  "credits_completed": 12
}
        \end{minted}
  \item Click \textbf{Send} and confirm you get a
        \texttt{201 Created} status back with the same object.
\end{enumerate}
 
\subsection{Why this matters once requests get complicated}
 
A single GET request is barely worth using Postman for; a browser
address bar does the job just as well. Postman starts to pay off once
your API involves things that are painful to keep retyping in a
terminal:
 
\begin{itemize}
  \item \textbf{Collections.} Every request you build can be saved into
        a named \newterm{collection}, so your entire API (every route in
        this chapter, for example) can be saved, organized, and rerun with
        a click instead of retyped.
  \item \textbf{Environments.} A variable like
        \texttt{\{\{base\_url\}\}} can be swapped between
        \texttt{http://localhost:3000} and a real deployed server without
        editing every saved request individually.
  \item \textbf{Authorization tab.} Instead of manually adding an
        \texttt{Authorization} header on every request, you can set your
        API key once per collection and Postman attaches it automatically.
  \item \textbf{Automated testing}: Postman can run a series of requests and check that the responses
        match expected values, so you can verify your API is working as
        intended after every change. This is great for testing before you have a front-end built, or for running regression tests after a code change.
\end{itemize}
 
Neither tool replaces the other. cURL is scriptable and works anywhere
you have a terminal, including inside automated tests and shell
scripts, which is why you'll keep seeing it later in this book. Postman
is better suited to exploring an API interactively, organizing a large
set of requests, and sharing that organized set with teammates. Most
professional API developers end up using both, for different parts of
their workflow.

\section{Summary}

This chapter touched on a lot of software and development concepts. Here are some of the key takeaways:
\begin{itemize}
  \item API's are ways to connect and fetch data from a server. They are a way to connect two different applications together.
  \item API's are usually built with a RESTful architecture, which is a set of principles that follow HTTP protocol to perform CRUD operations on resources.
  \item \textbf{CURL} is a command line tool that allows you to send HTTP requests to an API. Another commonly used tool is \textbf{Postman}, which is a GUI application that allows you to send HTTP requests to an API.
  \item We created a basic API that allows us to perform CRUD operations on a student resource. We can create, read, update, and delete students from our API.
\end{itemize}